\section{A proposed proof for all minimum cycle lengths at least seven}
\label{sec:all-cutoffs}

\textbf{Proposed theorem.}
For every pair of integers $n\ge k\ge7$, the polynomials $P_n=E_{n,k}/t$
and $P_{n+1}=E_{n+1,k}/t$ have respectively $n-2$ and $n-1$ simple,
strictly negative roots, and they interlace strictly.
The larger row supplies the two outer roots.

The rows themselves occur as the associated Stirling Eulerian polynomials
$A_n^{k-1}$ of Rakotoson and Rakotondrajao
(\href{https://cs.uwaterloo.ca/journals/JIS/VOL28/Rakotoson/rakoto12.html}
{\emph{J. Integer Seq.} 28 (2025), Article 25.5.7, Section~5}).
Their Theorems 30--32 give the insertion, distinguished-cycle, and generating
function identities. The conjecture addresses the zeros of these rows.

For the parameter decomposition put $x=n/k$. The insertion argument
already gives the theorem for $x<4$. At $x\ge4$ the two common-pole
comparisons treat all old sizes $n\ge10{,}000$ (with overlapping
aspect intervals $[4,48]$ and $[44,\infty)$). Through $n=250$, the
remaining 7,021 old-size and threshold pairs in the finite-prefix
certificate include every required pair.
For $251\le n<10{,}000$ and $x\ge20$, use the phase
$\Psi=(n+1)\theta-\Im G_s(s+i\pi)$ and the actual consecutive normalized
rows $Y_n,Y_{n+1}$. In the angular range $n\sin^2\theta\ge1$ the
bounds $|Y_n-\cos\Psi|<1/20$ and
$|Y_{n+1}-\cos\theta\,Y_n+\sin\theta\sin\Psi|<3\sin\theta/25$
give the next sign at an old zero. First hits of the phase levels give
ordered gaps, while the exterior estimate places the innermost ranks
before the first required hit. The precise contour, constants and
first-hit comparison are in the
\href{https://github.com/ProgramComputer/open-math-research/tree/research/topics/long-cycle-excedances/aspect-20}{large-aspect proof}.

It remains to consider $4\le x<20$. Use the parametrization
$t=-\exp(-\pi\cot\theta)$, $0<\theta\le\pi/2$, and set
$H=\lfloor7\sqrt n/20\rfloor$. At $n\ge1000$, the bounded-aspect
exterior estimate places $H$ simple roots of $P_n$ before the integer
phase sample at $(H+1)\pi/n$, giving the sign $(-1)^h$ of the
actual $P_{n+1}$ at each of these ranks $h$. The grouped pole estimates
give alternating signs at consecutive integer samples and the prescribed
next sign at each intervening old zero when $n\sin^2\theta\ge6/5$.
This last condition holds at the sample $(H+1)\pi/n$ and persists to the
midpoint. These two estimates therefore join without an unchecked gap.
The joint value and derivative bounds and their finite rational
majorants are in the
\href{https://github.com/ProgramComputer/open-math-research/tree/research/topics/long-cycle-excedances/bounded-aspect}{bounded-aspect proof}.

For $251\le n\le1000$, the finite exterior computation again supplies
these $H$ ranks with their signs. Put $\theta_j=j\pi/n$ and
$y_j=\exp(-\pi\cot\theta_j)$. Use the entrance threshold $A=7$
at $x\le15/2$, and $A=5$ at $15/2<x\le10$. At $10<x<20$ take
$A=5$ for $k\le100$ and $A=6$ for $k>100$. Choose an integer
$J\ge H+1$ with $k\sin\theta_J\ge A$ and
$n\sin^2\theta_J\ge6/5$. The high-angle estimate supplies the old and
actual next signs from that sample to the midpoint.

Here is the finite assertion needed at each $H+1\le j<J$: there
are $y_j\le a<b\le y_{j+1}$ with
\[
 (-1)^{j-1}P_n(-a)>0,\qquad (-1)^j P_n(-b)>0,\qquad
 (-1)^j P_{n+1}(-y)>0\quad(a\le y\le b).
\]
It follows for all the required gaps from the accompanying
\href{https://github.com/ProgramComputer/open-math-research/tree/research/topics/long-cycle-excedances/intermediate}{finite sign certificate}.
The coefficients used in the certificate are integer prefixes of the literal
cycle enumerators. An omitted term in the reduced polynomial has
absolute value at most $p^N y^{p-1}$ in row $N$. Indeed, cutting the
permutation into $p$ decreasing cyclic runs and labeling those runs gives
$p!$ encodings per permutation, among at most $p^N p!$ encodings.
For an upper bound $B\ge y_{j+1}$ the omitted tail after $W$ terms
is therefore bounded by
\[
 \frac{(W+1)^N B^W}{1-B((W+2)/(W+1))^N}.
\]
The denominator is checked positive. The ratio of successive majorants
decreases with $p$, which proves this geometric bound. The first program
evaluates $W=64$ terms with outward-rounded interval arithmetic on the
entire next-row interval. The second uses $W=60$ and proves the specified
next sign at both ends. In the latter method, the new signs at the adjacent
endpoints of consecutive old-root brackets force zeros in their
complementary gaps. Exhaustion of the new degree suffices for strict
interlacing. The two checks each cover 93,902 pairs, of which 82,395
require at least one intervening gap, and 541,199 such gaps in total.

In increasing modulus the first half-axis has
$H+(J-H-1)+(\lfloor n/2\rfloor-J)=\lfloor n/2\rfloor-1$
selected old roots. Reflection supplies their reciprocals. In odd size the
root $-1$ accounts for the remaining old degree and the high-angle
comparison gives its next sign. Thus exactly $n-2$ old roots have been
located, proving simplicity. The signs at zero and at the far negative
end complete the $n-1$ complementary gaps of the next polynomial, counting
both outer gaps. This concludes the proposed proof for $k\ge7$.

The all-size result at threshold three is Theorem~1. Bowen Zuo's supplying
releases separately state all-size interlacing for $k=4,5,6$, starting
with the old sizes $4,5,6$, respectively. Those are external inputs; their
pinned proof audits are recorded
\href{https://github.com/ProgramComputer/open-math-research/pull/216}{separately}.
There is no end-to-end Lean or Rocq formalization of the argument for $k\ge7$.
